% 第 52 章骨架：导引框已定，正文待写。写作计划见 docs/PLAN.md“第 52 章”，写法见 docs/STYLE.md。
\chapter{恶性通胀怎样开始、怎样结束}\label{ch:hyperinflation}

\begin{chapterguide}
\guidequestion{为什么有的国家一夜之间止住了恶性通胀，有的却反复失败？}
\begin{guidegoals}
\item 推导卡甘（1956）货币需求模型与恶性通胀的动态；
\item 用萨金特（1982）的“四次大通胀的终结”说明财政体制转变的作用；
\item 比较玻利维亚（1985）、以色列（1985）、巴西雷亚尔计划（1994）与阿根廷的多次稳定计划；
\item 用匈牙利（1946）、津巴布韦与委内瑞拉检验这些结论。
\end{guidegoals}
\guideprereq{\chref{ch:weimar}\enspace\nameref{ch:weimar}；\chref{ch:budget}\enspace\nameref{ch:budget}；\chref{ch:ratex}\enspace\nameref{ch:ratex}}
\end{chapterguide}

\chapterstub
